\chapter{Memoria descriptiva}

\section{Generalidades}

\subsection{Descripcion}

La presente memoria descriptiva corresponde al proyecto de instalaciones electricas interiores de una vivienda unifamiliar de dos niveles, ubicada en el distrito de San Miguel, provincia de San Roman, departamento de Puno. El documento se formula como parte del desarrollo academico del curso de Instalaciones Electricas y toma como base la informacion levantada para la vivienda de Aquiles Taylor Ramos Yapo.

El proyecto comprende la descripcion tecnica general del sistema electrico interior, desde el punto de alimentacion de baja tension hasta los circuitos derivados que atenderan los ambientes de la vivienda. La memoria se apoya en los croquis arquitectonicos disponibles, en las respuestas del cuestionario del usuario y en el trabajo CAD actualmente desarrollado para regularizar la distribucion de los pisos.

Los calculos de maxima demanda, seleccion final de conductores, protecciones, caida de tension, metrados y especificaciones completas se desarrollaran en los capitulos correspondientes. En esta memoria solo se establecen los criterios descriptivos y el alcance tecnico de la instalacion.

\subsection{Nombre del proyecto}

La presente memoria descriptiva corresponde al proyecto:

\begin{quote}
\textbf{Instalaciones Electricas Interiores de Vivienda Unifamiliar -- Caso Aquiles Taylor Ramos Yapo.}
\end{quote}

El nombre final podra ajustarse segun el formato exigido por el curso o por la revision academica.

\subsection{Ubicacion}

De acuerdo con las fuentes directas revisadas, la vivienda se ubica en la zona urbana del distrito de San Miguel. La referencia consignada para el predio es Avenida Horacio con jiron Marineros, manzana F7, lotes 11 y 12.

\begin{table}[H]
\centering
\caption{Ubicacion referencial del proyecto}
\begin{tabularx}{0.82\textwidth}{>{\bfseries}p{4cm}L}
\toprule
Region / Departamento & Puno \\
Provincia & San Roman \\
Distrito & San Miguel \\
Referencia urbana & Avenida Horacio con Jr. Marineros, manzana F7, lotes 11 y 12 \\
Zona & Urbana \\
\bottomrule
\end{tabularx}
\end{table}

\textbf{Dato por confirmar:} direccion final normalizada, lote aplicable al expediente, coordenadas del predio y referencia catastral definitiva.

\subsubsection{Condiciones climatologicas}

La vivienda se emplaza en una zona altoandina del departamento de Puno, por lo que el diseno electrico debera considerar condiciones ambientales propias de la region: temperatura baja, temporada de lluvias, posible humedad en exteriores y exposicion de elementos ubicados fuera de la vivienda.

Para esta fase no se fija una altitud ni una tabla climatologica definitiva, porque no se ha validado una fuente especifica para el punto exacto del predio. En el desarrollo final se debera confirmar la informacion climatologica con una fuente tecnica o institucional, especialmente si se especifican equipos, canalizaciones exteriores o elementos de puesta a tierra.

\subsubsection{Datos de coordenadas}

\textbf{Dato por confirmar:} coordenadas geograficas del predio. Las capturas catastrales disponibles sirven como referencia visual, pero no reemplazan una coordenada verificada para el expediente.

\subsection{Objetivos}

\subsubsection{Objetivo general}

Describir el sistema de instalaciones electricas interiores requerido para atender la demanda residencial de la vivienda, considerando criterios de seguridad, funcionalidad, continuidad del servicio, mantenimiento y cumplimiento de la normativa tecnica aplicable.

\subsubsection{Objetivos especificos}

\begin{itemize}
  \item Definir el alcance general de la instalacion electrica interior de la vivienda.
  \item Identificar los ambientes y cargas principales considerados en el proyecto.
  \item Establecer la organizacion preliminar del tablero general y de los circuitos derivados.
  \item Señalar los criterios generales para alumbrado, tomacorrientes, canalizaciones y puesta a tierra.
  \item Precisar los datos que deben confirmarse antes de desarrollar calculos, especificaciones, metrados y planos electricos definitivos.
\end{itemize}

\subsection{Alcances del proyecto}

El proyecto comprende el diseno de las instalaciones electricas interiores en baja tension para una vivienda unifamiliar de dos pisos. El alcance se desarrolla sobre la base de los croquis arquitectonicos disponibles y de la interpretacion CAD vigente en la carpeta de trabajo del caso Aquiles.

Los componentes considerados son:

\begin{itemize}
  \item Alimentacion electrica desde el punto de suministro o medicion hasta el tablero general.
  \item Tablero general de distribucion de la vivienda.
  \item Circuitos de alumbrado interior.
  \item Circuitos de tomacorrientes de uso general.
  \item Circuitos para cocina, zona de servicio, lavadora, bomba de agua u otras cargas que los calculos justifiquen.
  \item Canalizaciones, cajas, accesorios, interruptores, tomacorrientes y salidas electricas.
  \item Sistema de puesta a tierra de proteccion.
  \item Planos electricos, diagrama unifilar, cuadro de cargas y detalles en fases posteriores.
\end{itemize}

No se considera en esta fase la aprobacion de planos electricos, el metrado final, la conversion a DWG, el presupuesto, la ejecucion de obra ni la validacion profesional para construccion.

\subsection{Normas de aplicacion}

El proyecto se formula tomando como referencia la normativa tecnica aplicable a instalaciones electricas interiores de vivienda. Las normas identificadas en el repositorio y consideradas para esta memoria son:

\begin{itemize}
  \item Codigo Nacional de Electricidad -- Utilizacion \cite{cne_utilizacion}.
  \item Norma Tecnica EM.010 Instalaciones Electricas Interiores del Reglamento Nacional de Edificaciones \cite{rne_em010}.
  \item Criterios de la empresa concesionaria Electro Puno S.A.A. para suministro residencial, sujetos a confirmacion documental.
  \item Normas graficas y simbologia electrica que se definiran en el capitulo de planos.
\end{itemize}

Los articulos, numerales y reglas especificas deberan citarse en los capitulos de calculo, especificaciones y planos una vez verificados en los archivos normativos del repositorio.

\section{Descripcion del proyecto}

\subsection{Descripcion arquitectonica de referencia}

La vivienda corresponde a una edificacion unifamiliar de dos niveles, de material predominante ladrillo y concreto, destinada a siete miembros. La distribucion arquitectonica aun se encuentra en regularizacion, por lo que los ambientes descritos en esta memoria tienen caracter referencial.

En las respuestas del cuestionario se registro la siguiente distribucion preliminar:

\begin{itemize}
  \item Primer piso: dos cuartos grandes, cocina pequena, pasadizo, escalera o plataforma intermedia y bomba de agua exterior asociada a pozo.
  \item Segundo piso: sala, dos cuartos con cama, cocina amplia o zona de servicio, cuarto de uso varios y bano.
\end{itemize}

El trabajo CAD mas reciente identifica con mayor claridad la geometria del segundo piso, donde se observa sala, cocina grande, oficina o ambiente de usos varios, cuarto con cama, pasadizo, bano, escalera y un dormitorio inferior con forma en L. Esta condicion debe mantenerse en los planos electricos, porque influye en la ubicacion de luminarias, interruptores, tomacorrientes y recorridos de canalizacion.

\textbf{Dato por confirmar:} compatibilizacion final entre la nomenclatura del cuestionario y los nombres de ambientes que aparecen en el plano CAD.

\subsection{Alimentador principal}

El alimentador principal suministrara energia desde el punto de medicion o entrega de la concesionaria hasta el tablero general de la vivienda. De manera preliminar, el sistema se considera monofasico, 220 V, 60 Hz, con suministro de Electro Puno S.A.A.

La ruta, tipo de acometida, ubicacion del medidor, calibre del alimentador y proteccion principal no se definen en esta memoria. Estos datos deberan calcularse y representarse en el diagrama unifilar y en los planos electricos.

\textbf{Dato por confirmar:} punto exacto de suministro, potencia contratada si existiera, condicion del medidor, tipo de acometida y recorrido del alimentador principal.

\subsection{Tablero electrico general}

El tablero electrico general sera el punto de distribucion de los circuitos derivados de la vivienda. Se recomienda ubicarlo en el primer piso, en una zona seca, accesible, visible para operacion y mantenimiento, y alejada de condiciones de humedad directa.

El tablero debera contener la proteccion general y las protecciones de cada circuito derivado. La capacidad nominal, numero de polos, reserva de espacios, barras, interruptores y esquema definitivo se estableceran en el capitulo de calculos justificativos.

\textbf{Dato por confirmar:} ubicacion exacta del tablero general en el plano arquitectonico.

\subsection{Alimentadores secundarios}

En caso de que el proyecto requiera subtableros o alimentaciones independientes hacia zonas especificas, estos se definiran como alimentadores secundarios. Para la vivienda actual no se adopta todavia un subtablero como solucion definitiva; se evaluara si el tablero general puede atender directamente todos los circuitos derivados o si conviene sectorizar por piso.

\subsubsection{Canalizaciones}

Las canalizaciones interiores podran ser empotradas, superficiales o mixtas, segun el estado constructivo de la vivienda, los recorridos disponibles y la compatibilizacion con muros, losas y acabados. En ambientes interiores se priorizara una instalacion ordenada, protegida y registrable mediante cajas adecuadas.

En exteriores, zonas de humedad o recorridos expuestos, la canalizacion debera seleccionarse con criterio de proteccion mecanica y ambiental. La especificacion final del tipo de tuberia y accesorios corresponde al capitulo de especificaciones tecnicas.

\textbf{Dato por confirmar:} si la instalacion sera nueva empotrada, superficial sobre vivienda existente o mixta.

\subsubsection{Cables alimentadores}

Los cables alimentadores y conductores de circuitos se seleccionaran por capacidad de corriente, caida de tension, temperatura, tipo de canalizacion, longitud de recorrido y proteccion asociada. En esta fase no se asignan secciones definitivas.

Los conductores deberan ser de cobre, con aislamiento adecuado para instalaciones interiores de baja tension y con identificacion de fase, neutro y conductor de proteccion segun corresponda.

\subsubsection{Uniones y empalmes de conductores}

Las uniones y empalmes deberan ejecutarse solamente en cajas accesibles, con elementos adecuados para garantizar continuidad mecanica, continuidad electrica y aislamiento equivalente al conductor. No se deberan dejar empalmes ocultos dentro de tuberias o tramos no registrables.

\subsubsection{Uniones y empalmes dentro de canalizaciones}

No se admitiran empalmes dentro de canalizaciones cerradas. Todo empalme requerido debera resolverse en caja de paso, caja de salida, tablero o punto de registro diseñado para tal fin.

\subsection{Tableros electricos de distribucion}

El proyecto considera inicialmente un tablero general para toda la vivienda. Si en la fase de calculos se justifica una distribucion por pisos o por zonas, podran incorporarse subtableros, siempre que queden representados en el diagrama unifilar y en los planos.

\subsubsection{Interruptores termomagneticos}

Los circuitos deberan protegerse mediante interruptores termomagneticos seleccionados de acuerdo con la corriente de diseño, capacidad de los conductores y condiciones de cortocircuito esperadas. La capacidad nominal de cada interruptor se definira en los calculos justificativos.

\subsubsection{Interruptores diferenciales}

Se considerara proteccion diferencial en los circuitos que correspondan por seguridad de personas, especialmente tomacorrientes, zonas humedas y circuitos de uso residencial. La sensibilidad y ubicacion de estas protecciones se definiran en el esquema del tablero.

\subsection{Circuitos derivados}

Los circuitos derivados atenderan las salidas de alumbrado, tomacorrientes y cargas especificas de la vivienda. Como criterio preliminar, se evaluaran circuitos separados para:

\begin{itemize}
  \item Alumbrado del primer piso.
  \item Tomacorrientes generales del primer piso.
  \item Cocina pequena del primer piso, si su uso lo requiere.
  \item Alumbrado del segundo piso.
  \item Tomacorrientes generales del segundo piso.
  \item Cocina amplia o zona de servicio del segundo piso.
  \item Lavadora.
  \item Bomba de agua exterior asociada al pozo.
\end{itemize}

Esta sectorizacion proviene de las respuestas preliminares del cuestionario y sera revisada en los calculos justificativos. No constituye aun el numero definitivo de circuitos.

\subsubsection{Tubos de plastico PVC}

Las tuberias de PVC se consideraran para recorridos interiores protegidos, siempre que sean compatibles con el tipo de instalacion y las condiciones de servicio. El diametro, clase y accesorios se definiran por cantidad de conductores, seccion, facilidad de cableado y normativa aplicable.

\subsubsection{Tubos metalicos}

Las tuberias metalicas podran emplearse en tramos expuestos, zonas que requieran mayor proteccion mecanica o recorridos exteriores, si la solucion constructiva lo justifica. Su uso definitivo dependera de la compatibilizacion con arquitectura y condiciones de montaje.

\subsubsection{Soportes y accesorios}

Los soportes, abrazaderas, conectores, curvas, uniones y demas accesorios deberan asegurar continuidad, fijacion y proteccion de las canalizaciones. En instalaciones superficiales, los soportes deberan ubicarse de manera ordenada y permitir inspeccion visual.

\subsubsection{Cajas para salidas}

Las cajas para salidas de alumbrado, interruptores y tomacorrientes deberan instalarse en ubicaciones accesibles y coherentes con el uso del ambiente. Su tipo y dimension se definiran segun el numero de conductores, dispositivos y canalizaciones que confluyan en cada punto.

\subsubsection{Cajas de paso}

Se usaran cajas de paso cuando el recorrido, el cambio de direccion, la longitud de canalizacion o la cantidad de conductores lo requiera. Estas cajas deberan quedar registrables y no ocultas por acabados permanentes.

\subsubsection{Cinta aislante y elementos de aislamiento}

Los materiales de aislamiento usados en empalmes y terminaciones deberan ser compatibles con la tension de servicio y con el tipo de conductor. Su especificacion final se desarrollara en el capitulo de materiales.

\subsection{Iluminacion general}

El sistema de iluminacion debera cubrir dormitorios, sala, cocina, bano, pasadizos, escalera y zonas de servicio. Las respuestas del cuestionario registran puntos de alumbrado preliminares por ambiente, pero la cantidad final se definira luego de ajustar el plano arquitectonico.

\subsubsection{Nivel de iluminacion}

Los niveles de iluminacion se determinaran segun el uso de cada ambiente y la normativa aplicable. En esta fase no se asignan valores definitivos de iluminancia, porque corresponde verificarlos en el desarrollo de calculos y planos.

\subsubsection{Tipos de lamparas utilizadas}

Se considerara el uso de luminarias LED por eficiencia energetica, disponibilidad y bajo consumo. El tipo de luminaria, potencia, temperatura de color, ubicacion y forma de montaje se definiran en planos y especificaciones.

\subsection{Tomacorrientes}

Los tomacorrientes se distribuiran segun el uso de cada ambiente. En dormitorios y sala se consideraran puntos de uso general para equipos comunes. En la cocina amplia del segundo piso se deberan evaluar cargas como microondas, lavadora, waflera y otros artefactos de consumo medio o alto.

El bano cuenta con un tomacorriente preliminar, el cual debera ubicarse con criterio de seguridad y proteccion adecuada frente a humedad. No se considera ducha electrica; la informacion disponible indica uso de terma solar.

\textbf{Dato por confirmar:} cantidad final de tomacorrientes en sala, dormitorios, cuarto de usos varios y cocina.

\subsection{Sistema de puesta a tierra}

La vivienda no cuenta con un sistema de puesta a tierra confirmado. Por seguridad, el proyecto debera incluir puesta a tierra de proteccion para el tablero, tomacorrientes con contacto de tierra y partes metalicas que correspondan.

La ubicacion del pozo, tipo de electrodo, conductor de puesta a tierra, caja de registro y resistencia objetivo se definiran en el capitulo de calculos justificativos y en los detalles de planos.

\subsection{Colores en los conductores}

La identificacion de conductores debera seguir el criterio normativo aplicable para fase, neutro y conductor de proteccion. En el sistema monofasico previsto, se debera diferenciar claramente el conductor activo, el neutro y el conductor de puesta a tierra.

La codificacion final de colores se precisara en las especificaciones tecnicas y en la leyenda de planos.

\subsection{Pruebas electricas}

Antes de considerar operativa la instalacion, deberan realizarse pruebas de continuidad, aislamiento, verificacion de polaridad, accionamiento de protecciones y revision del tablero. Estas pruebas se documentaran en el expediente final o en los anexos tecnicos, segun el alcance del curso.

\subsubsection{Pruebas de aislamiento}

Las pruebas de aislamiento deberan realizarse por circuito y por alimentador antes de conectar cargas definitivas. Los valores admisibles y el procedimiento se definiran segun la normativa aplicable y las recomendaciones del equipo de medicion.

\subsection{Simbolos}

Los simbolos electricos empleados en planos deberan ser uniformes y reconocibles para alumbrado, interruptores, tomacorrientes, tablero, puesta a tierra, canalizaciones y cargas especiales. La simbologia se desarrollara en el capitulo de planos.

\subsection{Cuadro de cargas y maxima demanda}

El cuadro de cargas y la maxima demanda se desarrollaran en el capitulo de calculos justificativos. Para esta memoria solo se reconoce que las cargas principales preliminares son alumbrado, tomacorrientes generales, cocina, lavadora, bomba de agua exterior y equipos residenciales comunes.

No se consigna una potencia instalada ni maxima demanda definitiva en esta fase, porque las areas, cantidades finales de puntos y potencias reales de equipos aun deben confirmarse.
